\section{Discussion}
\label{sec:disc}

\subsection{What the evidence state buys, and where}
The stratified results support a specific reading of the method. GAST does not improve forecasting by being a larger or better-tuned model; its advantage is concentrated exactly where the observation record is degraded, and it vanishes or reverses where the record is complete or nearly absent. Three facts form this argument. The relative gain over M0 grows with severity across the high-missing range ($-3.9\%$ at five missing frames and $-7.3\%$ at six, against $+1.0\%$ when nothing is missing). The terminal-missing group, which is the condition that the anchor machinery targets, shows a 3.6\% reduction while the non-terminal group does not. And the anchor-lag curve is flatter than M0's over lags one to six. The gain is not strictly monotone in the missing count---the middle buckets fluctuate within a percentage point---but its concentration in the buckets where evidence is scarce is what a capacity explanation does not predict.

The reversal at lag seven is equally informative. When the anchor is eight frames old, the gap-aware velocity is estimated from a single surviving observation pair and is dominated by noise, while the reliability gate still prefers that stale state over the aggregate. A model that instead leans on a learned motion prior can be more accurate. This suggests that the evidence weighting should saturate: beyond some age, the reliability of a state should be constant rather than continuing to decay, and the decoder should be allowed to trust the prior. We did not tune this boundary on test data, and we report it as a limitation of the current formulation rather than as a hyperparameter choice.

\subsection{Scene heterogeneity}
The per-scene results are not uniform: GAST is the best method on ETH-M, UNIV-M and ZARA2-M, is behind M0 on ZARA1-M and HOTEL-M, and is behind AgentFormer on HOTEL-M by a wide margin. We can offer a hypothesis, although the benchmark does not permit a clean test of it. HOTEL-M and ZARA1-M are scenes with relatively slow, regular motion and strong scene constraints, where a unimodal extrapolation from the last observation is already accurate; the endpoint-conditioned decoder, which is designed to place probability mass on distinct destinations, adds variance without adding accuracy. ETH-M and UNIV-M contain more heterogeneous motion and more frequent direction changes, where the destination is genuinely uncertain and an explicit endpoint hypothesis helps. If this reading is correct, the decoder should be conditioned on a scene- or sample-level estimate of endpoint uncertainty rather than applied uniformly. We regard this as the most promising direction for closing the scene-level gap, but the present evidence is correlational.

\subsection{Coverage and ranking}
The dissociation between best-of-20 coverage and top-1 ranking is the most consequential negative result in this study. GAST improves the metrics that the benchmark ranks by, and degrades the metric that a downstream planner most often consumes, by roughly 11--13\%. Since the encoder-only variants do not show the degradation, the decoder is the plausible locus, and the diversity term plus the trajectory-level classification target are the plausible mechanisms. A direct test would be to retrain with the diversity weight annealed to zero and with a joint classification target, and to measure both metric families; we did not run this ablation, and we therefore state the diagnosis as a hypothesis. Until it is resolved, GAST should be used where best-of-$K$ coverage matters, and its top-1 mode should not be treated as a calibrated intention estimate.

\subsection{Limitations}
\textbf{Single seed.} All numbers are from one seed per scene. We have no confidence intervals, and we observed a 4.4\% run-to-run difference on UNIV-M between two runs of the same configuration during development. Differences of a few percent on a single scene, such as the 1.7\% macro margin over M0, are therefore within the range that a second seed could plausibly move. The macro ordering against the external baselines is far outside that range, but the ordering among the internal variants is not. A multi-seed study with paired, trajectory-clustered bootstrap is required before any claim of a stable improvement over M0.

\textbf{Protocol dependence.} TrajGap-Bench is a controlled stress test built on synthetic masks. Its missingness distribution is not claimed to match any real tracker, and the benchmark is restricted to five ETH/UCY scenes with an 8-frame history. External validity on real perception output and on other datasets (e.g.\ SDD or Argoverse) is untested.

\textbf{Test-set dependence.} The test split contains the same latent trajectory under multiple masks. Our aggregations treat samples as independent, which overstates the effective sample size for any significance statement. The released per-sample diagnostics include the source track identifier precisely so that clustered inference can be performed later.

\textbf{Ablation granularity.} As discussed in Section~\ref{sec:ablation}, the build-up sequence changes encoder and decoder together, so the contributions of GERE, ECTA and the goal decoder are not separated. The UT-ENC arm was also produced under the earlier checkpoint monitor.

\textbf{Missing comparisons.} We do not report a two-stage imputation baseline inside the direct table, because the setting exposes the model to more information and must be reported separately; the comparison is therefore absent rather than negative. We also do not report inference latency, which matters for deployment.

\subsection{Implications for benchmark design}
Two design choices in the protocol proved to be load-bearing. Merging Easy and Hard into a single Mixed split prevents a model from being selected on the difficulty at which it happens to be strongest, and it makes the missing-count stratification a diagnostic rather than a leaderboard. Fixing the checkpoint to the final head's validation $\mathrm{minFDE}_{20}$ removes a degree of freedom that otherwise silently differs across papers. Both choices are cheap to adopt and, in our experience, change conclusions: several variants that looked better under a per-difficulty protocol did not survive the Mixed contract. We recommend that future work on missing-history forecasting report at least the missing-count and terminal-missingness stratifications, since the aggregate can hide a reversal of the kind documented at seven missing frames.
